\documentclass[11pt,a4paper]{ctexart}
\usepackage[margin=2.2cm]{geometry}
\usepackage{xcolor}
\usepackage{enumitem}
\usepackage{amssymb}
\usepackage{booktabs}
\usepackage[normalem]{ulem}
\usepackage{array}
\usepackage{longtable}
\usepackage{titlesec}
\usepackage{hyperref}
\definecolor{navy}{HTML}{1A365D}
\hypersetup{colorlinks=true,linkcolor=navy,urlcolor=navy}
\setlist{nosep}
\titleformat{\section}{\Large\bfseries\color{navy}}{}{0pt}{}
\titleformat{\subsection}{\large\bfseries\color{navy}}{}{0pt}{}
\pagestyle{plain}
\renewcommand{\arraystretch}{1.25}

\title{\textbf{九年级上册语法总结}\\[4pt]
\large 人教版（新课纲版）九年级上册 · 按教材书末 Grammar 附录整理}
\author{初中英语资料站}
\date{2026-11-12 · 附配套专项练习 192 题}
\begin{document}
\maketitle
\section*{使用说明}
\begin{itemize}[leftmargin=1.4em]
\item 本讲义按九年级上册教材书末 Grammar 附录整理，表格与例句取自教材原文。
\item 网页版：\url{https://1901151458-prog.github.io/english-materials/g9-grammar/}
\end{itemize}

\section{教材 Grammar 附录总览}
\noindent\textcolor{gray}{九上教材书末的 Grammar 附录共三大板块，本站按原顺序整理。}\par\smallskip
人教版九年级上册教材在书末设有 \textbf{Grammar} 附录，用三张表格把全册语法收拢。原文结构如下：

{\footnotesize
\begin{longtable}{p{4.47cm}p{4.47cm}p{4.47cm}}
\hline
板块 & 教材标题 & 核心内容 \\
\hline
\endfirsthead
\hline
板块 & 教材标题 & 核心内容 \\
\hline
\endhead
\hline
\endfoot
一 & 动词的语态（Voices of Verbs） & 主动语态与被动语态；被动语态的构成（\textbf{表1}）；被动语态的用法 \\
二 & 宾语从句（Object Clauses） & 常见关联词及其用法（\textbf{表2}）：that / whether·if / wh- 词 \\
三 & 关系从句（Relative Clauses） & 关系代词引导的限制性关系从句（\textbf{表3}）：that / which / who \\
\end{longtable}
}

\begin{center}\fbox{\parbox{0.92\textwidth}{\small \textbf{说明：}「it 句型」在教材 Grammar 附录中未单列，但属于九上重点句式（配合 Unit 2 Inspiring People 等话题）。本站把它作为\textbf{补充专题}单列一讲，并接入已有专项练习。}}\end{center}

\subsection*{三个板块一句话记住}

\begin{enumerate}[leftmargin=1.4em]
\item \textbf{动词的语态}：主语是动作的承受者时，用 \texttt{be + 过去分词}。
\item \textbf{宾语从句}：一个完整句子作动词的宾语，靠关联词引导，用陈述语序。
\item \textbf{关系从句}：一个句子修饰前面的名词（先行词），靠关系代词连接。
\end{enumerate}

\section{动词的语态（Voices of Verbs）}
\noindent\textcolor{gray}{主动语态 vs 被动语态；构成、变化与用法}\par\smallskip
\subsection*{1. 主动语态和被动语态}

英语动词有两种语态，即\textbf{主动语态}（active voice）和\textbf{被动语态}（passive voice）。当主语为动作的执行者时，谓语要用主动语态；当主语为动作的承受者时，谓语要用被动语态。

\begin{itemize}[leftmargin=1.4em]
\item Many people \underline{speak} English.\\ 许多人说英语。（主动语态：主语 many people 是动作 speak 的执行者。）
\item English \underline{is spoken} by many people.\\ 英语被许多人使用。（被动语态：主语 English 是动作 speak 的承受者。）
\end{itemize}

\begin{itemize}[leftmargin=1.4em]
\item Alexander Bell \underline{invented} the telephone in 1876.\\ 亚历山大·贝尔在 1876 年发明了电话。（主动语态）
\item The telephone \underline{was invented} by Alexander Bell in 1876.\\ 电话是亚历山大·贝尔在 1876 年发明的。（被动语态）
\end{itemize}

\subsection*{2. 被动语态的构成}

被动语态由「\textbf{助动词 be + 及物动词的 -ed 形式}」构成。助动词 be 有人称、数和时态的变化。以动词 \texttt{ask} 为例：

{\footnotesize
\begin{longtable}{p{3.35cm}p{3.35cm}p{3.35cm}p{3.35cm}}
\hline
时态 & 肯定式 & 否定式 & 疑问式 \\
\hline
\endfirsthead
\hline
时态 & 肯定式 & 否定式 & 疑问式 \\
\hline
\endhead
\hline
\endfoot
一般现在时 & I \textbf{am asked} …\\ He / She \textbf{is asked} …\\ We / You / They \textbf{are asked} … & I \textbf{am not asked} …\\ He / She \textbf{is not asked} …\\ We / You / They \textbf{are not asked} … & \textbf{Am I asked} …?\\ \textbf{Is he / she asked} …?\\ \textbf{Are we / you / they asked} …? \\
一般过去时 & I \textbf{was asked} …\\ He / She \textbf{was asked} …\\ We / You / They \textbf{were asked} … & I \textbf{was not asked} …\\ He / She \textbf{was not asked} …\\ We / You / They \textbf{were not asked} … & \textbf{Was I asked} …?\\ \textbf{Was he / she asked} …?\\ \textbf{Were we / you / they asked} …? \\
一般将来时 & I / He / She / We / You / They \textbf{will be asked} … & I / He / She / We / You / They \textbf{will not be asked} … & \textbf{Will I / he / she / we / you / they be asked} …? \\
\end{longtable}
}

\subsection*{3. 被动语态的用法}

\begin{itemize}[leftmargin=1.4em]
\item \textbf{不知道或没必要说明执行者}时用被动；
\item \textbf{需要特别强调动作承受者}时用被动。
\end{itemize}

\begin{itemize}[leftmargin=1.4em]
\item A mini-camera \underline{is connected} to the front of the glasses.\\ 一个小型摄像头被连接在眼镜前方。（不必说明是谁连的。）
\item Some of the old buildings \underline{were destroyed} in the earthquake.\\ 一些老建筑在地震中被毁。（强调承受者 the old buildings。）
\item More astronauts \underline{will be sent} to the space station.\\ 更多宇航员将被送往空间站。（将来时被动：will be + 过去分词。）
\item The moon \underline{will be visited} by tourists one day.\\ 月球总有一天会被游客造访。（需要时可用 by 引出执行者。）
\end{itemize}

\begin{center}\fbox{\parbox{0.92\textwidth}{\small \textbf{语篇与语体：}教材指出，从语篇方面看，出于上下文衔接连贯的需要也会使用被动语态；从语体方面看，被动语态在\textbf{科技或新闻报道等正式文体}中使用频率更高。}}\end{center}

\section{宾语从句（Object Clauses）}
\noindent\textcolor{gray}{一个完整句子作宾语，靠关联词引导}\par\smallskip
在复合句中，\textbf{用作宾语的句子}叫宾语从句。引导宾语从句的常见关联词有 \texttt{that}、\texttt{if}、\texttt{whether}、\texttt{what}、\texttt{who}、\texttt{where}、\texttt{why} 和 \texttt{how} 等。

{\footnotesize
\begin{longtable}{p{6.70cm}p{6.70cm}}
\hline
关联词 & 例句 \\
\hline
\endfirsthead
\hline
关联词 & 例句 \\
\hline
\endhead
\hline
\endfoot
\texttt{that}\\ （在口语或非正式文体中常省略） & Mary thinks \textbf{(that)} the teams were just fantastic.\\ I believe \textbf{(that)} reading is important for learning.\\ Grandma still remembers \textbf{(that)} you love cookies. \\
\texttt{whether}、\texttt{if}\\ （在口语中常用 if） & I wonder \textbf{if / whether} there will be races again next year.\\ I want to know \textbf{if / whether} you read English books outside class.\\ I wonder \textbf{if / whether} you could help me. \\
\texttt{who}、\texttt{what}、\texttt{which}、\texttt{when}、\texttt{where}、\texttt{how}、\texttt{why} & Can you tell me \textbf{who} she is?\\ He cannot remember \textbf{what} he ate for breakfast.\\ I can’t remember \textbf{which} book I borrowed last month.\\ Excuse me, do you know \textbf{when} the bookstore closes?\\ Do you know \textbf{where} he lives?\\ I want to know \textbf{how} memory works.\\ I don’t know \textbf{why} I am so forgetful. \\
\end{longtable}
}

\paragraph{语序}

\begin{itemize}[leftmargin=1.4em]
\item 从句一律用\textbf{陈述语序}：关联词 + 主语 + 谓语。
\item × Do you know where \sout{does he live}?\\ ✓ Do you know where \textbf{he lives}?
\end{itemize}

\paragraph{省略}

\begin{itemize}[leftmargin=1.4em]
\item \texttt{that} 在口语中可省略；
\item \texttt{if / whether} 和 wh- 词\textbf{不能}省略。
\end{itemize}

\subsection*{拓展 A：主句变成被动后，从句怎么办？}

常见的 \texttt{People say that …} 这类句子改成被动后有两种处理，要注意 \textbf{that 从句的身份变了}。

{\footnotesize
\begin{longtable}{p{6.70cm}p{6.70cm}}
\hline
主动句（that 从句是宾语从句） & 被动句（that 从句变主语从句） \\
\hline
\endfirsthead
\hline
主动句（that 从句是宾语从句） & 被动句（that 从句变主语从句） \\
\hline
\endhead
\hline
\endfoot
People \textbf{say} that he is a genius. & \textbf{It is said that} he is a genius. \\
People \textbf{believe} that the story is true. & \textbf{It is believed that} the story is true. \\
They \textbf{report} that the team won. & \textbf{It is reported that} the team won. \\
\end{longtable}
}

\begin{center}\fbox{\parbox{0.92\textwidth}{\small \textbf{关键区别：}主动句里 that 从句作 say 的宾语，是\textbf{宾语从句}；改成 \texttt{It is said that} 之后，it 是形式主语，that 从句升格为\textbf{主语从句}。两者不要混为一谈。}}\end{center}

还有一种被动句里，从句\textbf{仍然是宾语从句}——因为被动动词后面本身还带宾语：

\begin{itemize}[leftmargin=1.4em]
\item I \underline{was told} that he would come.\\ 我被告知他会来。（that 从句作 was told 的宾语。）
\item He \underline{was asked} whether he liked it.\\ 有人问他是否喜欢它。（whether 从句作 was asked 的宾语。）
\end{itemize}

\subsection*{拓展 B：whether / if 从句移到句首}

宾语从句里 whether 和 if 常可互换；但一旦从句\textbf{移到句首作主语}，就\textbf{只能用 whether}。

\begin{center}\fbox{\parbox{0.92\textwidth}{\small \textbf{为什么：}句首的 \texttt{If} 会被理解成「如果」，引导条件状语从句，句子就缺主语了。所以表示「是否」时，从句放句首一定要用 whether。\\ \textbf{对比：}\textbf{Whether} he will come is unknown.（是否 → 主语从句）／\textbf{If} he comes, I’ll tell you.（如果 → 条件状语从句）}}\end{center}

如果不想把从句放句首，用 it 作形式主语更自然：

It is not clear \underline{whether he will come}.\\ 他是否会来还不清楚。

\subsection*{拓展 C：宾语从句的简化}

当主句和从句的主语一致（或从句主语泛指）时，宾语从句可以缩成「\textbf{疑问词 + 不定式}」，句子更简洁。

{\footnotesize
\begin{longtable}{p{6.70cm}p{6.70cm}}
\hline
宾语从句 & 简化为「疑问词 + 不定式」 \\
\hline
\endfirsthead
\hline
宾语从句 & 简化为「疑问词 + 不定式」 \\
\hline
\endhead
\hline
\endfoot
I don’t know \textbf{what I should do}. & I don’t know \textbf{what to do}. \\
He didn’t know \textbf{how he could answer}. & He didn’t know \textbf{how to answer}. \\
She wondered \textbf{whether she should go}. & She wondered \textbf{whether to go}. \\
Can you tell me \textbf{where I can find} it? & Can you tell me \textbf{where to find} it? \\
\end{longtable}
}

\paragraph{可以简化}

\begin{itemize}[leftmargin=1.4em]
\item 主句主语 = 从句主语：I don’t know what \textbf{I} should do.
\item 从句主语泛指、不出现时。
\end{itemize}

\paragraph{不能简化}

\begin{itemize}[leftmargin=1.4em]
\item 主语不一致：I don’t know what \textbf{he} should do.（不能缩成 what to do）
\item \texttt{that} 引导的陈述从句不能这样简化。
\end{itemize}

\begin{center}\fbox{\parbox{0.92\textwidth}{\small \textbf{注意：}简化后不再有「主语 + 谓语」，语序问题也就不存在了——这也是 \texttt{what to do} 比 \texttt{what should I do} 更保险的原因。}}\end{center}

\begin{center}\fbox{\parbox{0.92\textwidth}{\small \textbf{配套专项练习：}宾语从句站内已有 6 讲精讲 + 48 道分级练习，覆盖引导词、语序、时态呼应与教材 Unit 3 / Unit 4 语境。\\ 进入宾语从句专项 · 48 题 →}}\end{center}

\section{关系从句（Relative Clauses）}
\noindent\textcolor{gray}{也叫「定语从句」，修饰前面的先行词}\par\smallskip
关系从句是由\textbf{关系词}（关系代词或关系副词）引出的修饰名词或代词的分句结构。因为这种分句的主要功能是作名词或代词的修饰语（即「定语」），所以关系从句亦称\textbf{「定语从句」}。关系从句通常置于它修饰的名词或代词之后，被修饰的名词或代词叫作\textbf{先行词}。

\begin{itemize}[leftmargin=1.4em]
\item I love traditional music \underline{which / that tells a story}.\\ 我喜欢讲述故事的传统音乐。（先行词 music；关系从句作定语。）
\item I prefer music \underline{which / that is full of joy and celebration}.\\ 我更喜欢充满欢乐与庆典气氛的音乐。（先行词 music。）
\end{itemize}

引导关系从句的词有\textbf{关系代词} that、which、who（宾格 whom，所有格 whose）和\textbf{关系副词} where、when、why。关系代词和关系副词在从句中充当特定的成分。教材此附录重点介绍由关系代词 \textbf{that、which、who} 引导的\textbf{限制性}关系从句。

{\footnotesize
\begin{longtable}{p{4.47cm}p{4.47cm}p{4.47cm}}
\hline
关系代词 & 指代 & 例句 \\
\hline
\endfirsthead
\hline
关系代词 & 指代 & 例句 \\
\hline
\endhead
\hline
\endfoot
\texttt{that} 在从句中作主语或宾语 & 指物或指人 & I prefer music \textbf{that} touches me.（作主语）\\ That’s the only thing \textbf{(that)} I can do now.（作宾语）\\ He is the singer \textbf{(that)} you are looking for.（作宾语） \\
\texttt{which} 在从句中作主语或宾语 & 指物 & The song \textbf{which} uses the violin and the erhu is my favourite.（作主语）\\ He is listening to the music \textbf{(which)} he downloaded just now.（作宾语） \\
\texttt{who} 在从句中分别作主语和宾语 & 指人 & He is the kind of man \textbf{who} never gives up.（作主语）\\ The man \textbf{(who)} we met at the library this morning is our new teacher.（作宾语） \\
\end{longtable}
}

\begin{center}\fbox{\parbox{0.92\textwidth}{\small \textbf{教材原注：}在非正式文体中，关系代词在从句中作宾语时常可省略。\\ 注意：省略只发生在\textbf{作宾语}时；作主语的关系代词不能省。}}\end{center}

\section{非限制性关系从句（逗号 + which）}
\noindent\textcolor{gray}{用逗号隔开的补充说明，不能用 that，也不能省略}\par\smallskip
教材附录强调「此处重点介绍由关系代词 that、which、who 引导的\textbf{限制性}关系从句」。但阅读和写作中还会遇到另一种：\textbf{非限制性关系从句}——用\textbf{逗号}与主句隔开，只作补充说明，不限定先行词的范围。删掉它，主句意思仍然完整。

\subsection*{两种常见用法}

\begin{itemize}[leftmargin=1.4em]
\item My brother\textbf{, who lives in Beijing,} is a doctor.\\ 我哥哥是医生，他住在北京。（who 补充说明先行词 my brother。）
\item The book\textbf{, which I bought last week,} is very interesting.\\ 这本书很有意思，是我上周买的。（which 补充说明先行词 the book。）
\item He passed the exam\textbf{, which made his parents happy}.\\ 他通过了考试，这让他的父母很高兴。（\textbf{which 指代整个主句}，不是某个词。）
\item She was late\textbf{, which was unusual}.\\ 她迟到了，这很不寻常。（which 指代「她迟到」这件事。）
\end{itemize}

\subsection*{非限制性从句的四条规则}

\begin{enumerate}[leftmargin=1.4em]
\item \textbf{必须用逗号}：从句前后各一个逗号（在句末时只留前面一个）。
\item \textbf{不能用 that}：指物用 \texttt{which}，指人用 \texttt{who / whom}，所有格用 \texttt{whose}。
\item \textbf{不能省略关系词}：与限制性从句「作宾语可省略」不同，非限制性一律不能省。
\item \textbf{指代整句时只能用 which}：此时 which 相当于「这件事」，不能用 that、what。
\end{enumerate}

{\footnotesize
\begin{longtable}{p{4.47cm}p{4.47cm}p{4.47cm}}
\hline
对比项 & 限制性关系从句 & 非限制性关系从句 \\
\hline
\endfirsthead
\hline
对比项 & 限制性关系从句 & 非限制性关系从句 \\
\hline
\endhead
\hline
\endfoot
逗号 & 无 & \textbf{必须有} \\
作用 & 限定先行词范围 & 补充说明，不限定 \\
能否删除 & 不能（删后意思变） & 能（主句仍成立） \\
关系词 & that / which / who & which / who / whom / whose（\textbf{不用 that}） \\
作宾语省略 & 可省略 & \textbf{不可省略} \\
先行词 & 名词或代词 & 名词、代词，\textbf{或整个主句} \\
\end{longtable}
}

\begin{center}\fbox{\parbox{0.92\textwidth}{\small \textbf{易错：}× He was late, \sout{that} made the teacher angry. → ✓ He was late, \textbf{which} made the teacher angry.\\ \textbf{易错：}× My father, \sout{that} works in a hospital, is a doctor. → ✓ My father, \textbf{who} works in a hospital, is a doctor.}}\end{center}

\section{it 句型专题（补充）}
\noindent\textcolor{gray}{It is + adj. + for / of sb + to do sth}\par\smallskip
本专题在教材 Grammar 附录中未单列，但属于九上高频句式（配合 Unit 2 Inspiring People 等话题），单独收一讲。

\subsection*{基本结构}

{\footnotesize
\begin{longtable}{p{4.47cm}p{4.47cm}p{4.47cm}}
\hline
句式 & 例句 & 说明 \\
\hline
\endfirsthead
\hline
句式 & 例句 & 说明 \\
\hline
\endhead
\hline
\endfoot
\texttt{It is + adj. + to do sth} & It is important \textbf{to learn English well}. & \textbf{it} 是形式主语，真正主语是后面的不定式。 \\
\texttt{It is + adj. + for sb + to do sth} & It is hard \textbf{for us} to finish the work today. & 形容词描述\textbf{动作}（hard / easy / important / necessary）用 \textbf{for}。 \\
\texttt{It is + adj. + of sb + to do sth} & It is kind \textbf{of you} to help me. & 形容词描述\textbf{人的品质}（kind / nice / clever / foolish）用 \textbf{of}。 \\
\end{longtable}
}

\subsection*{for 还是 of：一看形容词修饰谁}

\paragraph{用 for：形容词说「事」}

\begin{itemize}[leftmargin=1.4em]
\item It is important \textbf{for us} to learn English well.
\item It is difficult \textbf{for him} to get up early.
\item 常见词：important、necessary、hard、easy、difficult、useful
\end{itemize}

\paragraph{用 of：形容词说「人」}

\begin{itemize}[leftmargin=1.4em]
\item It is kind \textbf{of you} to help me.
\item It is clever \textbf{of him} to work out the problem.
\item 常见词：kind、nice、clever、polite、foolish、careless
\end{itemize}

\begin{center}\fbox{\parbox{0.92\textwidth}{\small \textbf{判断口诀：}把 sb 提到前面，能说「\textbf{sb is adj.}」→ 用 \textbf{of}（You are kind. ✓）；不能说「You are important.」时 → 用 \textbf{for}（You are important 意思变了，所以用 for）。}}\end{center}

\subsection*{同义句转换}

\begin{center}\fbox{\parbox{0.92\textwidth}{\small \textbf{配套专项练习：}it 句型站内已有 6 讲 + 48 道分级练习，覆盖 for / of 辨析、形式主语、同义句改写与 Unit 2 语境。\\ 进入 it 句型专项 · 48 题 →}}\end{center}

\section{专项练习入口}
\noindent\textcolor{gray}{两个专项各 6 级、48 题，支持在线作答与即时解析}\par\smallskip
\paragraph{被动语态专项 · 48 题}

\paragraph{宾语从句专项 · 48 题}

\paragraph{关系从句专项 · 48 题}

\paragraph{it 句型专项 · 48 题}

\begin{center}\fbox{\parbox{0.92\textwidth}{\small \textbf{建议顺序：}先按本页 01–06 讲通规则 → 再进入对应专项练习自测 → 错题回到小节复习。四个专项合计 \textbf{192 题}，覆盖九上语法全部考点。\\ \textbf{考前速查：}一页纸速查表 →（A4 单页，可打印）}}\end{center}

\end{document}
